\section{The compactification and its enhancement locus}\label{sec:compactification}

\subsection{Stable factors and the orthogonal embedding}
Fix a smooth projective K3 surface $S$ with $\operatorname{Pic}(S)=\ZZ H$,
where $H$ is ample. Let $g_0$ be its Ricci-flat K\"ahler metric in that
class, oriented so that primitive $(1,1)$ forms are anti-self-dual. Norms
and harmonic representatives without an explicit scale use $g_0$ and
unnormalized trace inner products. Write $V_0=\Vol(S,g_0)$.

Choose locally free $\mu_H$-stable bundles $A,B$ with
\begin{equation}\label{eq:factors}
 \begin{gathered}
 \operatorname{rk}A=\operatorname{rk}B=3,\qquad \det A=\det B=\OO_S,\\
 c_2(A)=7,\quad c_2(B)=10,\qquad A\not\simeq A^*,\quad B\not\simeq B^*.
 \end{gathered}
\end{equation}
Here existence includes the locally free, slope-stable and non-self-dual
conditions. We give the argument because these conditions determine both
the analytic application and the unbroken gauge algebra.

For $k=7,10$ the Mukai vector is $v_k=(3,0,3-k)$, with
$v_7^2=24$ and $v_{10}^2=42$. Both vectors are primitive. An equal
reduced-Hilbert-polynomial proper subsheaf would have first Chern class
zero and Euler characteristic divided by rank equal to $-1/3$ or
$-4/3$. Neither is possible for rank one or two. Thus $H$ is general
for these vectors. Yoshioka's nonemptiness and irreducibility theorem,
together with the slope-stability and local-freeness criteria in
\cite[Theorem~0.1, Proposition~0.5 and section~3.3]{yoshioka}, gives
nonempty slope-stable and locally free opens. In that notation the
integer $l$ equals three and $v_k^2\ge2l^2$. The exceptions to the
locally free criterion do not occur: the vectors have rank three,
whereas the spherical exception of the form
$(\operatorname{rk}v_0)v_0-\omega$ would have square rank.
Irreducibility makes the two opens intersect. Since $\operatorname{Pic}^0(S)=0$,
their determinant is trivial.

Dualization acts holomorphically as an involution on this smooth open.
At a fixed point $F$, simplicity makes an isomorphism $F\to F^*$
symmetric or skew-symmetric after scalar normalization. Odd rank
excludes the latter. The invariant tangent space is
$H^1(S,\mathfrak{so}(F))\simeq H^1(S,F)$: use
$\Lambda^2F^*\simeq F$ and the symmetric form. Stability and Serre
duality give $H^0(F)=H^2(F)=0$, and Riemann--Roch gives
\begin{equation}\label{eq:selfdual-dim}
 h^1(F)=k-6,\qquad h^1(\End F)=6k-16.
\end{equation}
The fixed locus of a holomorphic involution is smooth with its invariant
tangent space. One can see this directly by replacing local coordinates
$z$ centered at a fixed point by
$\tfrac12(z+D\iota_0\,z\circ\iota)$; their derivative is the
identity and the action becomes linear. The fixed locus therefore
cannot fill the stable locally free open. This proves that the choices
in \eqref{eq:factors} exist. Their distinct second Chern classes also
exclude isomorphisms between $A$ or $A^*$ and $B$ or $B^*$.

At enhancement set $E_0=A\otimes\CC^2$ and define the real gauge
bundle
\begin{equation}\label{eq:gauge-bundle}
 \mathcal V=(E_0\oplus B)_{\RR}\oplus\RR^{14},\qquad
 \mathcal V\otimes\CC=E_0\oplus E_0^*\oplus B\oplus B^*\oplus\OO^{14}.
\end{equation}
It has real rank 32, $w_2(\mathcal V)=0$, and
$p_1(\mathcal V)=-2c_2(E_0\oplus B)=-48=p_1(TS)$.
The weighted instanton number is
\begin{equation}\label{eq:instanton}
 2c_2(A)+c_2(B)=2\cdot7+10=24.
\end{equation}
There is consequently no five-brane source in the integrated Bianchi
identity \cite{honecker}. The Donaldson--Uhlenbeck--Yau correspondence
\cite{donaldson,uy} supplies determinant-flat HE connections, which are
ASD on $S$. They give the leading background with eight supercharges,
or six-dimensional $\mathcal N=(1,0)$ supersymmetry. Equation
\eqref{eq:instanton} is an integrated condition; it does not set the
gauge and tangent curvature densities equal pointwise. Their difference
sources the first correction in section~\ref{sec:coupled}.

The parallel orthogonal adjoint sections have compact algebra
\begin{equation}\label{eq:centralizer}
 \mathfrak h_0=\mathfrak{so}(14)\oplus\mathfrak u(2)\oplus\mathfrak u(1),
 \qquad \dim_\RR\mathfrak h_0=96.
\end{equation}
Indeed the endomorphism blocks give $\End(\CC^2)$ and $\CC$ by
simplicity. The additional exterior-square blocks have no parallel
sections. For example,
\[
 \Lambda^2(A\otimes\CC^2)=
 (\Lambda^2 A\otimes\operatorname{Sym}^2\CC^2)
 \oplus(\operatorname{Sym}^2 A\otimes\Lambda^2\CC^2),
\]
and $H^0(A\otimes A)=\Hom(A^*,A)=0$ by stability and
non-self-duality. The other cross blocks vanish for the same reason.
There are no maps to the trivial fourteen-dimensional factor.
This also explains why replacing $A$ by a pseudoreal rank-two factor
would give a different centralizer. All central factors here have
$c_1=0$; the abelian flux mass matrix of \cite[section~3.2]{honecker}
vanishes. Protection of the nonconstant-profile vector calculation
requires the separate argument in section~\ref{sec:coupled}.

\subsection{Harmonic coordinates and actual bundles}\label{subsec:harmonic}
General K3 formality results \cite{budurzhang} identify quadratic
deformation germs. The quantitative metric application below needs
the stronger coordinate specification in the original harmonic chart.

The harmonic space $H^1(S,\End A)$ has complex dimension 26 and a
quaternionic structure. Choose thirteen orthonormal symplectic pairs,
normalized using the parallel holomorphic two-form. Then
\[
 H^1(S,\End E_0)=H^1(S,\End A)\otimes\End(\CC^2)
\]
has coordinates $(X_j,Y_j)_{j=1}^{13}$ with the flat real metric
\begin{equation}\label{eq:flatmetric}
 \|\dot x\|^2=\operatorname{Re}\sum_{j=1}^{13}
 \tr(\dot X_j^*\dot X_j+\dot Y_j^*\dot Y_j).
\end{equation}
The multiplicity group acts by simultaneous conjugation; its center
acts trivially. The complex and real moment equations are
\begin{equation}\label{eq:moments}
 \sum_j[X_j,Y_j]=0,\qquad
 \sum_j\bigl([X_j,X_j^*]+[Y_j,Y_j^*]\bigr)=0.
\end{equation}
Only their zero level is used. Let $\Cone$ be this hyperk\"ahler
quotient by the effective group $\PU(2)$, with metric $\gC$ on its
smooth-simple part.

We record why these are exact local deformation coordinates, rather
than merely the quadratic part of an unspecified obstruction. Put
$d=\dbar_{E_0}$ and denote harmonic projection and Green operators
by $H_q,G_q$. The convergent harmonic recursion of
\cite[Lemma~2.1]{wuebben15} is
\begin{equation}\label{eq:recursion}
 \mathfrak a(x)=x-d^*G_2(\mathfrak a(x)\wedge\mathfrak a(x)),\qquad
 \mathfrak k(x)=H_2(\mathfrak a(x)\wedge\mathfrak a(x)).
\end{equation}
On a Ricci-flat K3 its obstruction is exactly
\begin{equation}\label{eq:quadratic}
 \mathfrak k(x)=H_2(x\wedge x).
\end{equation}
For completeness, a harmonic $\End E_0$-valued $(0,2)$-form is
$\xi\bar\sigma$ with $\xi$ parallel. This follows from the formula
for $d^*(\xi\bar\sigma)$ and the Bochner identity on the
zero-mean-curvature endomorphism bundle. If $d^*u=0$ and
$\gamma=\phi\bar\sigma$, the coefficients of
$u\wedge d^*\gamma$ and $d^*\gamma\wedge u$ are, up to common
nonzero scalars,
$-\sum_j u_{\bar j}\nabla_j\phi$ and
$\sum_j(\nabla_j\phi)u_{\bar j}$. Pair either with a parallel
$\xi$ and integrate by parts, keeping the matrix order. The resulting
terms contain $\sum_j\nabla_j u_{\bar j}=-d^*u$ and vanish.
Consequently both products have zero harmonic projection.
Writing $\mathfrak a=x+\beta$ with $\beta$ coexact and
$d^*$-closed eliminates all terms containing $\beta$ in the
obstruction and proves \eqref{eq:quadratic}. This is precisely the
local argument in \cite[Lemma~3.1]{wuebben20}; it uses none of that
paper's later global rank or no-bubbling hypotheses.

Serre duality identifies \eqref{eq:quadratic} with the first equation
in \eqref{eq:moments}; a common nonzero symplectic factor and the
rank-dual identification do not change its zero set. Fixing the
determinant removes no harmonic directions because $H^1(S,\OO)=0$.
The analytic quotient construction \cite{buchdahl} and HE reconstruction
\cite[Theorem~1.1]{wuebben15} therefore identify the balanced quotient
with the actual nearby HE moduli. This identification $\Psi$ is
holomorphic in these same coordinates: the holomorphic classifying
map and the reconstructed HE map assign the same bundle class to
each harmonic deformation. No further nonlinear reparametrization
is made.

The quantitative input, specialized to this fixed background, is
\begin{equation}\label{eq:metriccomparison}
 (1-Cr^2)\gC\ \le\ \Psi^*\gKS\ \le\ (1+Cr^2)\gC,
 \qquad r=\|x\|,
\end{equation}
on genuine smooth stratum tangents. Here $\gKS$ is the real harmonic
Kodaira--Spencer metric in the unnormalized Dolbeault norm. The same
constant works when positive orbit scales have arbitrary ratios. It
may depend on the fixed metric, bundles, ranks, harmonic dimensions
and elliptic gaps. This does not assert uniformly sized smooth
coordinate charts through every lower stratum. On the smooth-simple
locus $\gKS$ is hyperk\"ahler \cite{itoh}.

Let $\sigma_j$ be the Pauli matrices with
$\tr(\sigma_i\sigma_j)=2\delta_{ij}$ and
$[\sigma_i,\sigma_j]=2i\epsilon_{ijk}\sigma_k$. Our family is
\begin{equation}\label{eq:pauli}
 X_1=a\sigma_3,\qquad X_2=b\sigma_1,\qquad
 X_j=0\ (j\ge3),\qquad Y_j=0,\qquad a,b\in\RR.
\end{equation}
It satisfies both moment equations since its nonzero matrices are
Hermitian. If $ab\ne0$, a line invariant under $\sigma_3$ is an
eigenline, and $\sigma_1$ exchanges the two eigenlines. The
representation is therefore simple. Reconstruction gives an actual
stable rank-six bundle $E_{a,b}$ with $c_2=14$ and trivial determinant.
The local simple moduli space has complex dimension
$104-2\cdot3=98$.

On a nonzero axis the representation is the sum of two distinct
characters, corresponding to two nearby stable rank-three factors.
Upper semicontinuity preserves the other Hom vanishings after
shrinking the neighborhood. The unbroken algebra is thus
$\mathfrak{so}(14)\oplus\mathfrak u(1)^2$ generically,
$\mathfrak{so}(14)\oplus\mathfrak u(1)^3$ on an axis, and
\eqref{eq:centralizer} at the origin. Their dimensions are
93, 94 and 96, respectively.
